\subsection{单生成群}

设 \(a \in Z\) ；因为 \(aZ\) 是 \(Z\) ，元素之间的关系 \(x, y\) 于 \(Z\) 它表明“存在 \(z \in Z\) 的集合，使得 \(x - y = az\) ”是一个等价关系，我们一致同意，一劳永逸地将其写为 \(x \equiv y \pmod{a}\) 或 \(x \equiv y(a)\) 并称之为模 \(a\) 的同余。将 \(a\) 中条件的扩张 \(-a\) 由此得到一个等价关系，因此可以假设 \(a \geqslant 0\) ；对于 \(a = 0\) , \(x \equiv y(0)\) 意味着 \(x = y\) ，因此只有当 \(a \neq 0\) 时才得到与相等不同的关系：因此在下文中我们将假设 \(a > 0\) 除非另有说明。

对于 \(a > 0\) ，商 \(\mathbf{Z}\) 由同余关系 \(x \equiv y(a)\) ，即群 \(\mathbf{Z}/a\mathbf{Z}\) ，被称为有理整数模 \(a\) .

命题16. 设 \(a\) 为整数 \(>0\) . 整数 \(r\) 的集合，使得 \(0 \leqslant r < a\) 构成该等价关系的一个代表系 \(x \equiv y \pmod{a}\) 在 \(\mathbf{Z}\) .

上定义了一个处处有定义的合成律。若 \(x\) 是一个整数 \(\geqslant 0\) 存在（集合论，III，§5，第6号）整数 \(q\) 和 \(r\) 的集合，使得 \(x = aq + r\) 和 \(0 \leqslant r < a\) 和 \(x \equiv r \pmod{a}\) 的一个正规子群。若 \(x\) 是一个整数 \(\leqslant 0\) ，整数 \(-x\) 是 \(\geqslant 0\) 且由上述存在一个整数 \(r\) 的集合，使得 \(0 \leqslant r < a\) 和 \(-x \equiv r \pmod{a}\) 。写 \(r' = 0\) 中正规，若 \(r = 0\) 和 \(r' = a - r\) 中正规，若 \(r > 0\) ，则

\[
x \equiv - r \equiv r ^ {\prime} \pmod {a}
\]

和 \(0 \leqslant r' < a\) 。我们现在证明，如果 \(0 \leqslant r < r' < a\) ，则 \(r \not\equiv r' \pmod{a}\) 的一个子群。现在 \(r' - r < na\) 对于 \(n \geqslant 1\) 和 \(r' - r > na\) 对于 \(n \leqslant 0\) 的自同态，因此 \(r' - r \notin a\mathbf{Z}\) .

推论。设 \(a\) 为整数 \(> 0\) 。群 \(\mathbf{Z} / a\mathbf{Z}\) 有理整数模 \(a\) 是一个阶为 \(a\) .

命题17. 设 H 是 Z 的一个子群。存在唯一一个整数 \(a \geqslant 0\) 使得 H = aZ。

上定义了一个处处有定义的合成律。若 \(\mathbf{H} = \{0\}\) ，则 \(\mathbf{H} = 0\mathbf{Z}\) 的两个元素。假设 \(\mathbf{H} \neq \{0\}\) 子群 \(\mathbf{H}\) 有一个元素 \(x \neq 0\) 。则 \(x > 0\) 或 \(-x > 0\) ，因此 \(\mathbf{H}\) 有元素 \(>0\) 的一个映射。设 \(a\) 是最小元素 \(>0\) 在 \(\mathbf{H}\) 子群 \(a\mathbf{Z}\) 生成的 \(a\) 包含于 \(\mathbf{H}\) ；我们证明 \(\mathbf{H} \subset a\mathbf{Z}\) 的一个映射。设 \(y \in \mathbf{H}\) 。根据命题16，存在一个整数 \(r\) 的集合，使得 \(y \equiv r (\text{mod. } a)\) 和 \(0 \leqslant r < a\) 更不用说 \(y \equiv r (\text{mod. } H)\) 的自同态，因此 \(r \in \mathbf{H}\) 但这仅在以下情况下才可能 \(r = 0\) ，从而 \(y \in a\mathbf{Z}\) 的一个分解。整数 \(a\) 是唯一的：如果 \(\mathbf{H} = \{0\}\) ，那么必然 \(a = 0\) ，且若 \(\mathbf{H} \neq \{0\}\) ，整数 \(a\) 是……的阶 \(\mathbf{Z}/\mathbf{H}\) .

定义15. 如果一个群允许由单个元素组成的生成元系，则称其为单生成群。有限单生成群称为循环群。

每个单生成群都是交换群（第3节，命题2的推论2）。单生成群的每个商群都是单生成群（第3节，命题2的推论3）。

加法群 Z 是单生成的：它由 \(\{1\}\) 对于每个正整数 a，群 Z/aZ 是单生成的，因为它是 Z 的商群。

命题18. 一个有限阶的单一生成群 \(a\) （使用复合 \(\mathbf{Z}/a\mathbf{Z}\) 。一个无限单生成群同构于 \(\mathbf{Z}\) .

设 G 是一个单生成群（乘法记号），x 是 G 的一个生成元。恒等元 \(x^{m}x^{n} = x^{m+n}\) (§ 1，第 3 号，公式 (1))表明该映射 \(n \mapsto x^{n}\) 是从Z到G的同态。其像是G的一个包含x的子群，因此就是G。根据第5条定理3，群G同构于Z关于某个子群的商群，该子群必然具有aZ的形式（命题17）。若a \textgreater{} 0时，群G是阶为a的有限群；若a = 0，则群G同构于Z。

命题19。设 \(a\) 为整数 \(>0\) 的一个映射。设 \(\mathbf{H}\) 的一个子群，且包含于 \(\mathbf{Z}/a\mathbf{Z}\) , \(b\) 的顺序 \(\mathbf{H}\) 和 \(c\) 其在……中的索引 \(\mathbf{Z}/a\mathbf{Z}\) 。则 \(a = bc\) , \(\mathbf{H} = c\mathbf{Z}/a\mathbf{Z}\) 和 \(\mathbf{H}\) （使用复合 \(\mathbf{Z}/b\mathbf{Z}\) .

反之，设 \(b\) 和 \(c\) 是两个整数 \(>0\) 的集合，使得 \(a = bc\) 。则 \(a\mathbf{Z} \subset c\mathbf{Z}\) 和 \(c\mathbf{Z}/a\mathbf{Z}\) 是 \(\mathbf{Z}/a\mathbf{Z}\) ，其阶为 \(b\) 且指数为 \(c\) .

由第4号命题的推论4，a = bc。根据第7号定理4，H具有形式 \(H'/aZ\) 的x和y，其中 \(H'\) 是Z的一个子群，且 \(Z/H'\) （使用复合 \((\mathbf{Z}/a\mathbf{Z})/H\) 因此阶为 c。由命题 17 及命题 16 的推论， \(H' = cZ\) 因此 H 是单生成的。最后，由命题 18，H 同构于 Z/bZ。反之，若 a = bc，则 \(aZ \subset cZ\) 对于 \(a \in cZ\) ：商群 \((\mathbf{Z}/a\mathbf{Z})/(c\mathbf{Z}/a\mathbf{Z})\) 同构于 Z/cZ（第 7 节，定理 4），因此阶为 c（第 4 节，命题 4 的推论），指数为 b（第 4 节，命题 4 的推论）。

推论。单生成群的每个子群都是单生成的。

设 \(a\) 和 \(b\) 是两个整数 \(\neq 0\) 。关系 \(b \in a\mathbf{Z}\) 也写作： \(b\) 是 \(a\) 的倍数，并且也 \(a\) 整除 \(b\) 或 \(a\) 是……的因子 \(b\) .

定义16. 一个整数 \(p > 0\) 被称为素数，如果 \(p \neq 1\) 且它没有因子 \(> 1\) 除了 \(p\) .

命题20。一个整数 \(p > 0\) 当且仅当该群 \(\mathbf{Z} / p\mathbf{Z}\) 是单群。

这由命题19得出。

推论。每个交换单群都是素数阶循环群。

设 G 是这样的一个群。则 \(G \neq \{e\}\) ；令 \(a \neq e\) 是G的一个元素。由a生成的子群是正规的，因为G是交换的，它并不缩减为 \(\{e\}\) 因此等于G。所以G是单生成的，从而同构于形如Z/pZ的群，其中p \textgreater{} 0，因为Z不是单群，且p必然是素数。

注记。阶为素数的有限群G必然是循环群。G除G本身和 \(\{e\}\) 外没有其他子群，因此它由每个非单位元素生成。 \(\neq e\) .

引理2。设a为整数 \textgreater0。通过将每个合成列与 \((\mathbf{H}_{i})_{0\leqslant i\leqslant n}\) 在群 Z/aZ 中，序列 \((s_{i})_{1\leqslant i\leqslant n}\) 的x和y，其中 \(s_{i}\) 是……的阶 \(H_{i-1}/H_{i}\) ，从而在 Z/aZ 的合成列与有限序列 \((s_{i})\) 整数 \textgreater0，使得 \(a=s_{1}\ldots s_{n}\) 之间建立一一对应。合成列 \((\mathbf{H}_{i})_{0\leqslant i\leqslant n}\) 是 Jordan-Hölder 列当且仅当 \(s_{i}\) 为素数。

上定义了一个处处有定义的合成律。若 \((\mathbf{H}_i)_{0\leqslant i\leqslant n}\) 是 \(\mathbf{Z} / a\mathbf{Z}\) ，由此，通过对 \(n\) 进行归纳，根据第4号，命题4，可知 \(a = \prod_{i=1}^{n} (\mathbf{H}_{i-1}:\mathbf{H}_i)\) .

反之，设 \((s_i)_{1 \leqslant i \leqslant n}\) 是一个整数序列 \(>0\) 的集合，使得 \(a = s_1 \ldots s_n\) 的一个正规子群。若 \((\mathbf{H}_i)_{0 \leqslant i \leqslant n}\) 是 \(\mathbf{Z} / a\mathbf{Z}\) 的集合，使得 \((\mathbf{H}_{i-1} : \mathbf{H}_i) = s_i\) 对于 \(1 \leqslant i \leqslant n\) ，那么必然 \(((\mathbf{Z} / a\mathbf{Z}) : \mathbf{H}_i) = \prod_{1 \leqslant j \leqslant i} s_j\) 正如通过对 \(i\) 的自同态，因此 \(\mathbf{H}_i = \left( \prod_{j=1}^{i} s_j \right) \mathbf{Z} / a\mathbf{Z}\) 进行归纳所见（命题19），这显示了所讨论映射的单射性。我们现在证明其满射性。写出 \(H_{i}=\left(\prod_{j=1}^{i}s_{j}\right)\mathbf{Z}/a\mathbf{Z}\) 对于 \(0\leqslant i\leqslant n\) ，得到Z/aZ的一个合成列，使得 \((\mathrm{H}_{i-1}:\mathrm{H}_{i})=s_{i}\) （命题19）。第二个断言由命题20及第7号，命题9得出。

设 \(\mathfrak{P}\) 表示素数的集合。

定理7（分解为素因数）。设 \(a\) 设为一个严格正整数。存在且仅存在一个族 \((v_{p}(a))_{p \in \mathfrak{P}}\) 整数 \(\geqslant 0\) 的集合 \(p \in \mathfrak{P}\) 使得 \(v_{p}(a) \neq 0\) 是有限的且

\[
a = \prod_ {p \in \mathfrak {P}} p ^ {v _ {p} (a)}.\tag{5}
\]

当群 \(\mathbf{Z}/a\mathbf{Z}\) 是有限的，则它有一个若尔当-赫尔德序列。引理2随后蕴含 \(a\) 是素数整数的乘积，因此该族的存在性 \((v_{p}(a))\) ；进一步地，对于每一个族 \((v_{p}(a))_{p\in\mathfrak{P}}\) 满足定理7的条件，整数 \(v_{p}(a)\) 是，对于所有 \(p\in\mathfrak{P}\) ，等于一个Jordan-Hölder序列的因子数 \(\mathbf{Z}/a\mathbf{Z}\) 同构于 \(\mathbf{Z}/p\mathbf{Z}\) （引理2）。该族的唯一性 \((v_{p}(a))\) 因此由Jordan-Hölder定理（第7号，定理6）得出。

推论。设 \(a\) 和 \(b\) 是两个整数 \(>0\) 。则 \(v_{p}(ab) = v_{p}(a) + v_{p}(b)\) 。对于 \(a\) 整除 \(b\) ，必要且充分的条件是 \(v_{p}(a) \leqslant v_{p}(b)\) 对于每个素数 \(p\) .

在任意群 G 中，如果由元素 \(x \in G\) 生成的（单生成）子群的阶为有限数 d，则称 x 为 d 阶元素；因此数 d 是满足 \textgreater0，使得 \(x^{d} = e\) 的最小整数；如果由 x 生成的子群是无限的，则称 x 为无限阶元素。这些定义连同命题 4（第 4 节）特别蕴含：在有限群 G 中，G 的每个元素的阶都是 G 的阶的因子。

命题21。设G是一个阶为n的有限群， \(x^{n} = e\) 为所有 \(x \in G\) .

上定义了一个处处有定义的合成律。若 \(p\) 是……的阶 \(x\) ，则 \(n = pq\) ，其中 \(q\) 一个整数，因此 \(x^n = (x^p)^q = e\) .

